\chapter{Software-managed caching and overlays}\label{ch:swmanaged}

\IS has no hardware caches. When a program's access pattern is regular it streams data with
transfers; when it is irregular, the program builds a cache in its D-LS. Two hardware assists make the
common operations cheap without taking policy away from software: a tag lookup for data caches and an
overlay map for code.

\section{Software data caches}\label{sec:swcache}

\subsection{Geometry}
A software cache is described by three CSRs:

\begin{center}\small
\begin{tabular}{@{}l l@{}}
\toprule
\textbf{CSR} & \textbf{Contents}\\\midrule
\csr{ttag} & D-LS address of the tag table (8-byte aligned)\\
\csr{tdata} & D-LS address of the line storage ($\textit{sets} \times \textit{ways}$ lines)\\
\csr{tgeom} & bits 5:0 $\log_2$(line bytes); bits 11:6 $\log_2$(sets); bits 15:12 ways (1--15)\\
\bottomrule
\end{tabular}
\end{center}

\begin{figure}[htbp]
\centering
\begin{bytefield}[bitwidth=0.75cm,endianness=big,boxformatting={\centering\footnotesize}]{16}
\bitheader{0,5,6,11,12,15}\\
\bitbox{4}{ways}\bitbox{6}{$\log_2$ sets}\bitbox{6}{$\log_2$ line}
\end{bytefield}
\caption{Low 16 bits of \csr{tgeom}; higher bits are ignored}\label{fig:tgeom}
\end{figure}

The tag table holds, for each set, \textit{ways} tag doublewords followed by one replacement-counter
doubleword, so set $s$ starts at $\csr{ttag} + s \cdot 8(\textit{ways}+1)$. A tag doubleword holds
$(\textit{line} \ll 1) \mid 1$, where $\textit{line} = \textit{offset} \gg \log_2(\text{line bytes})$
is the full line number; zero means invalid. Software initialises the table (normally to zero), so the
format is architectural.

\subsection{Lookup and install}
\insref{ltag} \texttt{rd, rs} looks up far offset \reg{rs}:
\begin{itemize}
\item \emph{hit}: \reg{rd} = D-LS address of the cached byte, i.e.\ line storage plus the offset within
  the line;
\item \emph{miss}: \reg{rd} = D-LS address of the victim line's storage, with bit 63 set. The victim is
  the way selected by the set's replacement counter.
\end{itemize}
\insref{linst} \texttt{rs} writes the tag for \reg{rs}'s line into the victim way and advances the
counter, which cycles $0, 1, \ldots, \textit{ways}-1, 0, \ldots$ (round robin).

A miss is handled by ordinary code. Because of the as-if-synchronous rule the load that follows the
fill is always correct; there is no window in which the cache can return stale bytes:

\begin{asm}
load_cached:                            ; a0 = far byte offset -> a0 = value
        ltag  t0, a0                    ; look the offset up in the tag table
        bltz  t0, lc_miss               ; bit 63 set means a miss
        ld    a0, 0(t0)                 ; hit: t0 is the LS address of the word
        ret
lc_miss:
        slli  t0, t0, 1                 ; clear bit 63 of the victim slot address
        srli  t0, t0, 1
        srli  t1, a0, LINE_LG           ; t1 = line number
        slli  t1, t1, LINE_LG           ; t1 = far offset of the line start
        li    t2, 1 << LINE_LG          ; one line, in bytes
        tget  tg0, t0, c0, t1, t2       ; fetch the line into the victim slot
        linst a0                        ; record the new line in the tag table
        andi  t1, a0, (1 << LINE_LG) - 1 ; byte offset within the line
        add   t0, t0, t1                ; LS address of the wanted word
        ld    a0, 0(t0)                 ; stalls until the line has landed
        ret
\end{asm}

The \texttt{swcache} example (Section~\ref{sec:ex-swcache}) is a complete set-associative cache.
Write-back, write-allocate, prefetching and replacement policy are all decided by software; the
hardware supplies only the lookup.

\section{Code overlays and code paging}\label{sec:overlays}

Code that does not fit in the I-LS, or that is loaded on demand, lives in far memory and is copied in
with \insn{tgetx}. The \term{overlay map} lets a program call such code by its far address and have the
hardware take the fast path when it is already resident. It has $\textit{NXMAP}$ entries, each mapping
a far address in some space to an I-LS address. $\textit{NXMAP}$ is a property of the platform, 4 to 16
(16 by default), given by the platform table's header (Section~\ref{sec:stpd}).

\begin{description}[leftmargin=1em,style=nextline]
\item[\insref{xcall} \texttt{rd, cs, ro}] looks up far address $\texttt{C[cs].base} + \reg{ro}$ in the
  overlay map. On a hit it jumps to the mapped I-LS address and leaves a link in \reg{rd}. On a miss it
  traps with cause \cause{xmiss} and \csr{tval} = the far address.
\item[\insref{xret} \texttt{rs}] returns through a link left by \insn{xcall}.
\item[\insref{xmap} \texttt{idx, rl, cs, ro}] sets entry \field{idx} (0 to $\textit{NXMAP}-1$; a
  larger index is an illegal instruction) to map that far address to I-LS address \reg{rl}.
\item[\insref{xunmap} \texttt{idx}] clears entry \field{idx}.
\end{description}

The miss handler chooses where to put the code, loads it with \insn{tgetx}, records the mapping with
\insn{xmap} and returns with \insn{eret}, which re-executes the \insn{xcall}; this time it hits. The
handler does not wait for the code: fetching from the range being filled stalls. \insn{tgetx}
automatically invalidates every overlay entry whose I-LS slot overlaps the range it overwrites, so the
map never points at code that has been replaced.

Overlay code must be position-independent. All \IS branches and jumps are PC-relative, so code
assembled at one address runs at any other. The \texttt{overlay} example (Section~\ref{sec:ex-overlay})
runs three routines through two I-LS slots.

\subsection{Exact and slot mapping}\label{sec:xslot}

CSR \csr{xslot} selects how an entry matches:

\begin{description}[leftmargin=1em,style=nextline]
\item[\csr{xslot} = 0 (exact mapping)] An entry matches one far address. Each routine is called at its
  entry point and mapped individually. A link is always an ordinary I-LS return address.
\item[\csr{xslot} = $n > 0$ (slot mapping)] Code is organised in \term{units} of $2^n$ bytes, and the
  I-LS holds a cache of $2^n$-byte \term{slots}. A unit's far address must be a multiple of $2^n$:
  \insn{xmap} traps with \cause{xalign} (\csr{tval} = the offset \reg{ro}) otherwise. An entry maps the
  unit starting at its far address to the slot starting at its I-LS address: any far address inside
  the unit hits, at the same offset in the slot. A program can therefore call any function of a unit,
  not only its first, and a unit may contain many functions.
\end{description}

So a lookup compares the far address shifted right by \csr{xslot} with each entry's, and the target is
the entry's I-LS address plus the low \csr{xslot} bits of the far address. An entry made under a
smaller \csr{xslot} covers, under a larger one, the aligned slot that holds its far address: the low
bits of the entry are ignored. (Software that changes \csr{xslot} normally unmaps first.)

The bits of \csr{xslot} above bit 5 are ignored; it is zero at reset and writable in either mode.

\subsection{Far links}

With slot mapping, a unit may be evicted while code inside it is waiting for a call to return. An
\insn{xcall} executed inside a mapped slot therefore does not leave an I-LS return address: it leaves
a \term{far link}, the far address of the instruction after the call with bit~63 set. \insn{xret}
resolves a far link through the overlay map exactly as \insn{xcall} resolves a target. If the caller's
unit is still resident, \insn{xret} returns into it, wherever it now sits; if it has been evicted,
\insn{xret} traps with \cause{xmiss} and \csr{tval} = the far link, the handler pages the unit back
in, and \insn{eret} retries the return. An ordinary link (bit~63 clear) makes \insn{xret} behave like
\insn{jr}, so functions that may be called from either resident or paged code return with
\insn{xret ra}.

\subsection{Demand paging}

Slot mapping turns the overlay map into a demand-paging mechanism for code, with the policy entirely in
software. A miss handler:
\begin{enumerate}
\item takes the far address from \csr{tval}, clearing bit~63 if the miss was a return;
\item rounds it down to the start of its unit;
\item picks a victim slot (StrataOS uses round robin over the slots that fit between the end of the
  program's resident code and the end of the I-LS, at most one per overlay map entry);
\item loads the unit with \insn{tgetx}, which also unmaps whatever the slot held;
\item maps the unit with \insn{xmap} and returns with \insn{eret}.
\end{enumerate}
The \texttt{units} program of StrataOS (Section~\ref{sec:os-paging}) runs 24 one-KiB code units through
a code cache of sixteen slots (eight with an eight-entry overlay map) with call chains seven deep, and
checks its result against the same computation done by resident code.
